\documentclass[11pt]{amsart}
\usepackage[french]{babel}
\usepackage[T1]{fontenc}
\usepackage[utf8]{inputenc}
\usepackage{hyperref}

\begin{document}
\title{Revues}
\maketitle
\tableofcontents

\subsection*{\textsl{La domination masculine}, Pierre Bourdieu}$ $\\

\textbf{Résumé.} \textsl{La domination masculine}, publié en 1998, s'attache à dépasser la critique locale des relations de subordinations entre sexes et entreprend une \textit{déhistoricisation} des schèmes sous-jacents à ces dominations. Ainsi réactualise-t-il les institutions, non seulement produits mais causes des formes de domination auto-légitimante, comme process construits.\\

\textbf{Détail.} Bourdieu esquisse dans \textsl{Une image grossie} de rendre compte de la représentation qu'à prit la condition féminine dans l'ordre masculin occidental. Problème : celui qui en parle est lui aussi enchainé. Une analyse de la société des Berbères de Kabylie permet un double détachement : du corps et du langage. Cet ordre masculin tient sa principal force du fait qu'il \textit{passe pour} neutre, s'incorporant dans une distinction et justification anatomique uniquement. Cette ambiguïté est tout à fait illustré par la société Kabyle (le \textit{remplissement}, \textit{laâmara} signifie aussi le sperme ; ), et abouti une fois confondu par dominants et dominés à un structuralisme ordonné selon que « tous actes de \textit{connaissances}, sont inévitablement des actes de \textit{reconnaissance}, de soumission ». Ainsi les habitudes sexuelles et sociales sont différentiées par ce système symbolique – jusqu'à la limite de domination qu'est la \textit{possession}, simple transformation de la femme en part de capital symbolique et social –, véritable lieu de \textit{violence symbolique}, acte de connaissance et méconnaissance par la vision du dominant. (Comment émergent-elles ? Bourdieu relève des rites de séparation initiatiques clairs dans la société Kabyle.) 

Or ces structures de dominations sont non seulement perpétués par les agents dominés et dominants mais reproduite sous forme de signes aux seins des \emph{institutions} (cf. F. Lordon, \textsl{Vivre sans ?}, chapitre $1$ pour cette notion). Ainsi analyse-t-il d'abord dans l'\textsl{Anamnèse des constantes cachées} les moyens d'émancipations féminins par ceux-là même qui lui sont présentés comme tels. Le mariage, les gestes de politesses et la mode sont proposés à la femme comme élévation de son ordre, alors qu'il ne modifie sa vision que dans l'\textit{illusio}. Mais il fallait enfin « reconstruire l'histoire du travail historique de déshistoricisation » dans \textsl{Permanences et changements} à travers les récents changements qu'ont subit l'\'Eglise et l'\'Ecole. La première fonde une représentation patriarcale, et réduit l'exercice féminin à celui d'\textit{effort de pose}, de « propagande icconographique ». La seconde, en préparation de l'\'Etat, séparée de la première, contient en elle l'opposition primitive observée comme constante. Ainsi dans l'\'Etat s'opposent les structures financières aux structures sociales, et à l'Université les sciences (dures) aux humanités (molles). Cela malgré le fait que Bourdieu reconnait en l'accès des filles et femmes à l'enseignement supérieur un changement fondamental : non pas la solution, mais les capacités de solutions. Ces trois instances font ainsi \textit{violence} au lieu de domination interdit de toute influence \textit{a priori}, à savoir le domicile conjugal. 

Pour conclure, c'est donc en « tenant ensemble la totalité des lieux et des formes dans lesquels s'exerce cette domination que l'on peut saisir les constantes de sa structures et les mécanismes de sa reproduction », et aboutir non pas à un \textit{renversement}, mais à un \textit{périssement axiologique} de la domination masculine.
\newpage

\subsection*{\textsl{Surveiller et punir}, Michel Foucault}$ $\\

\textbf{Résumé.} Michel Foucault décrit dans \textsl{Surveiller et punir} (1975) l'évolution de, d'une part, l'évolution d'une justice criminelle à une justice punitive pour devenir au \textsc{xx}\ieme \ siècle \textit{discipline} pénitentiaire ; d'autre part les corps et âmes y sont étudiés, regardés et jugés : mise en lumière du jeu de l'auto-perception des corps sociaux ; « l'âme, effet et instrument d'une anatomie politique ; l'âme, prison du corps. »\\

\textbf{Détail.} La naissance de la prison est récente et nécessite une retrospective plus que centenaire. Ainsi \textsl{Supplice} commence-t-il par un supplice et un emploi du temps. Dès la suite des reformes pénales du \textsc{xviii}\ieme \ et \textsc{xix}\ieme \ siècles s'impose une nouvelle forme du jugement : on juge non seulement le crime commis, mais aussi l'âme du criminel. Ainsi la procédure punitive dépasse la seule punition physique, et atteint ce que Foucault nomme l'âme du condamné. Cette dernière « ne naît pas fautive ou punissable, mais naît plutôt de procédures de punition, de surveillance, de châtiment et de contrainte. Elle est l'élément où s'articulent les effets d'un certain type de pouvoir et la référence d'un savoir. » Punition du roi, souverain, incarnation de la loi meurtri par le criminel ; il s'agit donc de \textit{restaurer} sa souveraineté. En parallèle s'effectue une autre modification, celle de la perception populaire du supplice. Longtemps acceptée comme structurante et récréative, la punition publique où le corps est présenté comme purifié, puissant de subir ce supplice même, elle se déplace lentement vers un rejet d'une forme de violence atteignable. 

Ce changement marque la naissance de la \textsl{Punition}. Caractérisée par un adoucissement des pratiques punitives, d'un faible soutient populaire et \textit{a fortiori} d'une diminution des « crimes de sang » pour les délits, le châtiment se veut exemplaire. Foucault précise que la nouveauté de cette conception hérétique est son devenir « généralisable à travers tout le corps social, susceptible de coder tous les comportements. » L'allègement des pratiques s'accompagne donc naturellement d'un « pouvoir publicitaire » redoublé : bien plus que la peine, importe les \textit{signes} de la peine. \textsl{Surveiller et punir} appui alors sur la création de la détention comme cristallisation de cette affect législatif commun, dépassant les espérances initiales des réformateurs pénitenciers. 

Avant de s'attarder sur la naissance de la prison et de ses corollaires, \textsl{Discipline} revient sur le \textit{micro-pouvoir} qu'ont fait émerger le « quadrillage des corps » de l'instruction militaire et scolaire et scolastique. Le sujet, nécessairement borné en espace et en temps, devient un point dans la répartition disciplinaire ; et est réciproquement « l'unité n'y est ni le territoire (unité de domination), ni le lieu (unité de résidence), mais le \textit{rang} ». Suivant Tocqueville, cette isotropie du redressement (cf. \textsl{Panopticon}) se voit renforcer par normalisation, processus d'autant plus efficace qu'il est soutenu par l'\textit{illusio} – séparateur – de l'examen ; dont Foucault montre toute sa relativité. Ainsi le corps du supplicié n'est-il plus « le point d'application dressé contre le corps du roi ; mais bien l'individu disciplinaire. »

Mais alors s'opère un nouveau changement dans la fonction de la prison, qui n'est plus seulement lieu de représentation des peines mais établissement coercitif et avilissant. La détention (légale) s'accompagne non seulement d'un isolement, mais premièrement dans le fait même de rendre cet isolement \textit{utile} par une discipline déjà analysée des transformations locales ; et secondement d'un contrôle (visuel via l'organisation architecturale, mentale par les rapports médicaux, etc.) normalisateur. Par conséquent le détenu n'est plus hors-la-loi, mais dévie « insensiblement de la discipline à la loi, de la déviation à l'infraction. » Car pour Foucault le problème du carcéral se situe ici : le socle axiologique de la prison tient à des principes qui lui sont internes. Et naturellement la forme \textit{délinquance} fait suite à l'emprisonnement, solution dans la loi. Loi dont la stratégie n'est que rarement discuté.
\newpage

\subsection*{\textsl{Les héritiers}, Pierre Bourdieu et Jean-Claude Passeron}$ $\\

\textbf{Résumé.} L’École se voulant outil démocratique par essence, toute critique de ses pratiques ou du simple \textit{constat} de l'échec de celle-ci serait vaine. Mais n'est-ce pas légitimer le système éducatif, son fonctionnement et ses préférences ? Bourdieu et Passeron montre la perpétuation des inégalités sociales par la transmission \textit{héréditaire} des formes de savoir et de dire.\\

\textbf{Détail.} Il n'est d'autre lieu que l’École où se dévoile le mieux l'origine sociale, économique et donc culturelle des étudiants. L'accès à la connaissance nécessite préalablement le goût ; une connaissance préalable des \textit{choix} d'études ; une posture et des habitudes que les relevés de Bourdieu et Passeron montrent. Les enfants d'agriculteurs sont sur plusieurs plans (connaissances musicales, théâtrales, aisance orale, ressources et logement, etc.) souvent deux à trois fois moins préparés que les enfants de cadres moyens et supérieurs, ou de profession libérale. Ainsi les réussites personnelles sont souvent décorrélées d'un éventuel mérite (par exemple, les étudiants fils ou filles de paysans n'ayant pas eu de propédeutique à l’École avaient un niveau en langue latine supérieure à la moyenne. Ce phénomène – l'\textit{emballement préparatoire} – n'est pas propre à cette classe sociale, et est présent en toute inégalité de fond.). Ainsi le \textsl{choix des élus} est fait.

Quelle forme prend alors la communauté étudiante et son évolution ? Celle-ci possède une unité, fondée sur le \textsl{jeu du sérieux}. « S'il est vain de découvrir des modèles de comportement proprement étudiants sous des conformismes sporadiques et changeants, c'est peut-être que les étudiants se ressemblent surtout par la nature du rapport qu'ils entretiennent à ce qu'ils (font dire à ce qu'ils) font et ce qu'ils sont. » Ce que cela implique, en particulier,  est que les structures de culture et de pouvoir sont adoptées comme \textit{posture} face à soi : l'étudiant est d'abord celui qui apprend à reconnaître les valeurs de la culture — à l'inverse, la révolte contre la contrainte de la règle permet intériorisation des valeurs que cette dernière impose. L’École enferme dans une irréalité qui se joue réalité, il reste l'adhésion au jeu intellectuel qui est de nouveau une constante de classe.

En conclusion, le transfert réciproque s'opérant entre étudiants et étudiant-professeur (se) maintient (sur) l'idée de \textit{don} ; rétablir une équité culturelle sinon de réussite se sera pas sans une « rationalisation » de la pédagogie.
\newpage

\subsection*{\textsl{Raisons pratiques}, Pierre Bourdieu}$ $\\

\textbf{Résumé.} Par quels moyens un sujet est-il poussé à l'\textit{action}, et oriente ses \textit{choix} ? Rompant avec les perceptions sujetiste, motivationnelle, et les dualités qui traversent les espaces sociaux et symboliques, Bourdieu propose une philosophie de l'action \textit{relationnelle} et \textit{dispositionnelle}.\\

\textbf{Détail.} L'action est le résultat de disposition (c'est à dire d'habitus), et n'est par conséquent jamais corollaire d'une substance qui serait propre à un être ou a un groupe social (\textit{e.g.} les sports dit nobles ne le sont que parce que ces derniers connaissent et se reconnaissent dans ces pratiques). Plus précisément, c'est de la formation même de ces dispositions de corps, de temps et et d'emploi du temps que naisse les germes de trajectoires biographiques.

Les moyens de faire la différences sont nombreux, autant qu'il est possible d'en faire. Pour ce qui est de l’École, cf. \textsl{Les héritiers}. Bourdieu reprend dans \textsl{Espace social et espace symbolique} le principe de \textsl{La Distinction} : c'est en et par elle que s’institut les formes de discernement, de goût dit personnels, dont les contours sont les écarts et différences \textit{perçus}. Ainsi l'existence des classes sociales est-elle conditionnée à la distinction des classes d'habitus – rien d'autres. 

La cristallisation de ces habitus forment des \textit{champs} – social, scientifique, bureaucratique ou politique – dont la caractéristique d'appartenance primordiale repose sur l'accord tacite de l'\textit{illusio} : le jeu oublié en tant que jeu. Bourdieu analyse le champs artistique qui se distingue par son auto-censure et exclu le « fondement de l’indépendance à l'égard des conditions historiques, qui réside dans le processus historique qui a conduit à l'émergence d'un univers capable d'assurer à ceux qui l'habitent d'une telle indépendance. » Est-ce \textit{dénaturer} et rendre impure l’œuvre – ou la personne – que de la située, distinguée ? Les formes « d'universalités transhistoriques » ne s'instituent qu'au travers d'une production collective. Il n'y a de nature que sociale. Appliqué au champs scientifique, la règle de l'art est celle de la \textit{vérité}, et Bourdieu précise : « S'il y a vérité, c'est que la vérité est un enjeu de lutte. » 

Une fois la topographie axiologique établie, il s'agit de comprendre quelle est sa dynamique. Et d'abord le rôle de l’État dans la disposition symbolique de ces habitants : penser l’État, c'est encore une pensée d’État. Il vaut peut-être de rappeler que la majorité de nos structures et capacités de perceptions et distinctions sont bien le corrélat d'une construction lente et fragmentée ; malgré tout le \textit{naturel} dont elles se parent. Ce dernier se démarque des autres formes d'institution en ce qu'il a en plus du capital le \textit{méta-capital} de l'exercice du pouvoir. Le pouvoir de l’État s'est lui-même construit (Bourdieu revient sur le rôle de la noblesses et de ses intendants qui lui ont progressivement donner pouvoir de \textit{nommer}, de vêtir - via les \textit{lois somptuaires}, bref un pouvoir divin) et dire ce qu'il est est encore un moyen d'« abolir l'histoire », d’apparaître universel : « la domination au nom de l'universel pour accéder à la domination ».

Je finis avec deux citations fondamentales du \textsl{Point de vue scolastique}:

\hspace{0.3cm}

\begin{center}« Constat simple qui conduit à un programme éthique ou politique, lui-même très simple : on ne peut échapper à l'alternative du populisme et du conservatisme, deux formes d'essentialisme tendant à consacrer le \textit{statu quo}, qu'en travaillant à universaliser les conditions d'accès à l'universel. » 

\hspace*{0.2cm}

« Tout projet de développement de l'esprit humain qui, oubliant l'enracinement historique de la raison, compte sur la seule force de la raison et de la prédication rationnelle pour faire avancer les causes de la raison, et qui n'en appelle pas à la lutte politique pour tenter de doter la raison et la liberté des instruments proprement politiques qui sont la conditions de leur réalisation dans l'histoire, reste encore prisonnier de l'illusion scolastique. »\end{center}

\hspace*{0.4cm}

\begin{center}\textsc{Lexique} \end{center}

\textsl{Habitus} : Principe générateur – des principes de différentiation différents – et unificateur qui retraduit les caractéristiques intrinsèques et relationnelles d'une position en style de vie unitaire. 

\textsl{Champs} : Couple autonome de champs de forces et de luttes sur l'espace social (ou politique) dont les agents se trouvent structurés et structurants du second, s'affrontant selon leur (dis)positions dans le premier.

\newpage 

\subsection*{\textsl{La chauve-souris et le capital}, Andreas Malm}$ $\\

\textbf{Résumé.} La pandémie de 2020 introduit le monde entier à l'\textit{urgence chronique} depuis longtemps prospectée par les chercheurs en zoologie et bactériologie. Seulement les parallèles de cette crise à la crise climatique sont trop nombreux pour qu'ils soient décorrélés, c'est ce que montre Malm et point leurs causes.\\

\textbf{Détail.} Dans un premier temps, l'émergence conjointe de l'épidémie de Covid-19 et de nombreux méga-feux (Australie, un milliards d'animaux morts), migrations climatiques (22 millions d'êtres humains en 2020), événements graves et rares (invasion de criquets rasant la moitié des récoltes en Afrique), etc. montre que s'il y a un perdant, ce sont les pauvres. Ce serait-il passé la même prise de décision immédiate des gouvernements si les infections avaient eu lieu en Iran, Irak ou au Bangladesh ? (réponse : non, cf. virus Ebola, Nipah Zika, SRAS et MERS)

Le dénominateur commun à ces deux crises (dont aucune n'est ponctuelle ni sans effets à long terme) est le dérèglement climatique causé par un ensemble de facteurs précis (détaillés ci-dessous). Les \textit{débordements zoonotiques} (passages d'un élément pathogène d'une espèce à une autre) sont favorisés par la déforestation – deuxième cause de pollution par émission dans le monde –, l'érosion des sols qui lui est associé et les départs forcées d'espèces (les chiroptères et les rongeurs réunissent à eux seuls une majorité des pathogènes transmissibles tels que les coronavirus), la réduction de l'espace et du temps et l'universalisation de la marchandisation capitaliste (l'union européenne importe une surface équivalente à l'Inde de terre incorporée par la consommation par an). En résumé, le capital \textit{est} un parasite qui n'entre en relation avec le monde sauvage « qu'en s’attachant à lui de façon à lui faire relâcher des marchandises qui possède une valeur d'échange. »

Dans \textsl{communisme de guerre}, Malm analyse les perceptions guerrières qu'ont éprouvé les régimes communistes depuis la première guerre mondiale et ce qu'il faut en tirer. Les homologies de la révolution russe de 1917 prédite et annoncée par Lénine et la catastrophe écologique sont nombreuses : en particulier, la « sociale - démocratie », qui repose sur le principe du temps, ne convient plus. Ainsi une action immédiate consisterait à 1. réviser l'ensemble des chaînes de production et mettre fin aux exploitations de terres tropicales. 2. Reboiser les monocultures et paysages défrichés. 3. Nationaliser et planifier la production d'énergie et réquisitionner les « monstres » que sont Exxon, Engie, BP, China Coil, etc. \\

La solution ne sera ni effectuée à demi, ni créée avec les outils de production du capital. Seul l'auto-défense humaine et la prise de conscience de leur propre naturalité pourront lui permettre, avec Adorno, de « mettre fin à leur propre domination de la nature ».

\newpage 

\subsection*{\textsl{D'un siècle l'autre}, Régis Debray}$ $\\

\textbf{Résumé.} Entreprise introspective de Régis Debray à vocation \textit{aperçive}, sur les altérations qu'ont connu d'abord l'homme — dont les étapes seront l'\textit{\'Ecole,} la \textit{prison}, l'{\'Etat} et enfin le \textit{congé} —, puis son pays, et enfin l'Homme. Les mystères auxquels sa vie s'est accrochés sont ceux du politique, de la civilisation et de la religion ; une triade que l'ont retrouvera. Et du « cela aussi passera » s'ajoute « eadem sunt semper omnia » de Lucrèce. \\

\textbf{Détail.} D'abord l'\textsl{\'Ecole}, ayant pour fonction de plus en plus \textit{distincte} son propre dépassement : y apprendre à ne pas aimer l'école, est à l'image pour l'auteur de l'exergue prise au \textsl{Manifeste} de Marx et Engels : « Il en s'agit plus d'interpreter le monde mais de le transformer. » Y passe la transition du capitalisme au socialisme, à l'envers certes, sous l'égide philosophique et sociologique neuve alors portée par Sartre, Aron, Merleau-Ponty, Bergson, Althusser, etc. Debray préfère s'en référer à Hugo (« Ce n'est pas à l'appui d'un texte, en vérité, que l'on gravit les parois de la réalité. »), et résume alors sa trajectoire intellectuelle. D'abord mettre sa pensée en action. Ensuite, penser l'action. Enfin, étudier l'action de la pensée – médiologie.

Alors, la \textsl{Prison}. La remarque suivante \begin{center}« Tous les malheurs des hommes viennent de ce qu'ils ne savent pas se tenir en repos dans leur chambre. »\end{center} de Pascal en dit assez pour signifier la force de clarté induite. Trois ans en Bolivie après après accompagné le Che – \textit{anywhere} – et Castro – \textit{somewhere} – en 1967. Clarté sur les conditions des limites et colimites du \textit{nous}. S'il y a eu défense, c'est qu'il y a un \textit{ici} à défendre. Le territoire n'est même pas à ignorer, étant sens le \textit{contre}. Aller jusqu'à porter comme raison d'être l'altérité ? En tout cas Gramsci, « marxiste de la sortie du marxisme »,  parle sans honte de nation, et il faut conclure sur le lien du principe de communion-\textit{éthos} à celui de rébellion-\textit{démos}.

Cela va plus loin, il y a dans cet ennemi, loin de la victime et du bourreau des \textsl{Justes} de Camus, proche du juge et du porteur, « comme de la punition dans l'espace et de la promesse dans le temps. » L'institution est alors le seul mur pouvant définir l'ange et le militaire ; et repousser la gabegie et le fugace. Puisque la communauté nécessite d'abord du \textit{Faire}, l'anarchisme – lieu de l'\textit{Être} – n'a pu se réaliser que sous forme ponctuelle. Et elle devient paradoxalement que vecteur d'une « éthique pour consommation privée. » \textsl{Voir pour prévoir}. En congé, Debray remarquera la dimension surfacique du sens commun :  des articles de la Constitution de la République aux neuroleptiques langagiers, du \textit{nous} au \textit{moi-je}, de Auguste Comte aux contes augustes.

L'\textsl{\'Etat} donc, faut d'immanent mouvements communs du \textit{Faire}, pour dresser et détourner le regard de la Table des Lois, ou voir plus loin, Debray insiste (et sous des formes diverses, cf. \textit{supra}) sur la transcendance du liant. La vérité, politique ou mystique d'abord — dont Dostoïevski se vante bien de lui préférer le Christ « même démontré mathématiquement ». C'est en découvrant l'Asie, le Proche-Orient et l'Afrique que malgré lui la conclusion se veut biblique : « Nation, c'est narration », ou avec l'Ogre sur Robespierre, « La vérité sur l'histoire ne sera pas ce qui a eu lieu mais ce qui sera raconter. » S'il y a vérité, c'est par conséquent, il y a \textit{communication} – et Debray ajoutera par sa médiologie la \textit{transmission}. Ne pas \textit{reproduire} l'erreur nécessite le regard de la chouette de Minerve par laquelle \textsl{D'un siècle à l'autre} commence et termine. Quoi que, « Essayer. Rater. Essayer encore. Rater mieux. » conjointement à l'effet abrutissant de la réussite et stratégique de l'échec, le choix vaut la peine d'être pensé deux fois. Image wittgensteinienne de De Gaulle, dont \textsl{Le fil de l'épée} marque le caractère et sa \textit{praxis} du « refus du repos de l'esprit et du mystère de l'inconnu. » Le silence peut signifier – trivialité. Mieux : « Rien ne vaut rien. Il ne se passe rien et cependant tout arrive. Mais cela est indifférent. » (Nietzsche)

Point fondamental de la pensée de l'auteur, un théorème d'incomplétude analogue à ceux de Gödel, dont Debray avoue et nuance le douteux rapprochement. Le groupe, la communauté, se démarque par un mythe d'origine singulier. Ce qui fait exister le groupe est nécessairement extérieur à celui-ci. Conséquence troublante, le principe instituant une communauté est d'une autre nature que l'institué, et de fait, la société n'\textit{est pas} sociologique. À l'inverse un entre-soi sans fond ni plafond se disloque au premier coup venu. Pragmatique et déprimant, peut-on espéré une rationalité lucide comme seul fondement d'un groupe ? Les hommes ont certes vu évolué leur environnement matérielle, aucun progrès n'est noté dans les relations entre les hommes. « Y a-t-il un mieux dans l'amour ? Un mieux dans la grandeur d'âme, ou à l'inverse, dans la soif de puissance et le diviser pour régner ? Un mieux dans le crime passionnel ou l'appât du gain ? ». Debray prend acte de l'exergue de Lucrèce et admet la contingence de l'acte, du choix. 

Enfin la médiologie et l'exigence du temps de la transmission qui ne supplantera jamais les progrès techniques – réseaux sociaux, téléphones et instantanéité \textit{spatiale} partagée. Même l'\'Eglise fit l'erreur, le rapporte Alfred Loisy. Reste à clarifier l'usage du vocable de la croyance, toujours vu sous un \textit{culte}, une \textit{spiritualité} ou pire : \textit{religion} (latin : \textit{religar} ; lien fondateur, loin du rite...) Régis Debray conclut sur une question fondamentale : pour quoi sommes-nous prêt à mourir ? Quelles limites infranchissables décrétées ? 

\newpage 

\subsection*{\textsl{La reproduction}, Pierre Bourdieu et Jean-Claude Passeron}$ $\\

\textbf{Résumé.} Prolongeant l'analyse héréditaire des moyens et des pratiques de l'appropriation des savoirs par le système éducatif des \emph{Héritiers}, les sociologues montrent que l'action pédagogique dans sa forme la plus abstraite s'exerce par un arbitraire double. D'abord l'arbitraire du savoir transmis, et ensuite celui du pouvoir de transmission. La violence symbolique, déjà rencontrée, s'incarne en une violence légitime dans les représentations scolaires et docimologiques.

\textbf{Détail.} Cette dernière affirmation constitue en fait l'axiome préludant au livre I : 

\textsl{\quad 0. Tout pouvoir de violence symbolique – tout pouvoir qui parvient à imposer des significations et à les imposer comme légitimes en dissimulant les rapports de forces qui sont au fondement de sa force –, ajoute sa force propre, \emph{i.e.} proprement symbolique, à ces rapports de force.}

Axiome légitime en tant qu'illustration sociale de l'auto-sustentation de la puissance de la multitude par elle-même — et dont la méconnaissance ou l'érudition à son propos constitue son noyau (cf. ). Cependant, en marge d'une démarche spinoziste, les rapports entre les instances symboliques sont eux-même produits et reproduits par le pouvoir symbolique – selon l'\textsl{axiome} et la \textsl{proposition 4}. La force symbolique est intrinsèque à la place occupée dans le champ pédagogique.

\textsl{\quad 1. Toute \emph{action pédagogique (AP)} est objectivement une violence symbolique en tant qu'imposition par un pouvoir arbitraire, d'un arbitraire culturel.}

Premièrement, la culture et la signification propre à un groupe ou une classe ne peut se justifier que par l'axiomachie socio-logique. Il convient donc d'en évacuer toute prétention naturalisante ou essentialisante. Deuxièmement l'action pédagogique se transmet selon l'agencement du méta-champ institutionnel, et la reproduction s'effectue par et en elle. Ainsi les différences d'accès, de réussite dans l'enseignement supérier 

\textsl{\quad 2. En tant que pouvoir symbolique s'exerçant dans un rapport de communication qui ne peut produire son effet propre, l'AP implique nécessairement comme condition sociale d'exercice l'\emph{autorité pédagogique (AuP)} et l'\emph{autonomie relative} de l'instance chargée de l'exercer.}

On se voit plier à reconnaitre que l'effet induit en celui-ci un degré de légitimité. Mais cette légitimité de AuP est à double sens : celui qui prêche ne convainc que les convaincus.

\textsl{\quad 3. En tant qu'imposition arbitraire d'un arbitraire culturel supposant l'AuP, l'AP implique le \emph{travail pédagogique (TP)} comme travail d'inculcation formant un habitus capable de se perpétuer après cessation de l'AP.}

Ici encore les agents qui fournissent un travail pédagogique ne sont jamais mieux enfermés en pensée et en pratique que lorsque la discipline et la censure sont intériorisées, et qu'ils croient vivre dans « l'illusion de la liberté et de l'universalité. »

\textsl{\quad 4. Tout \emph{système d'enseignement institutionnalisé} doit les caractéristiques spécifiques de sa structure et de son fonctionnement au fait qu'il lui faut produire et reproduire, par les moyens propres de l'institution, les conditions de l'institution, de l'arbitraire culturel, et des rapports entre les groupes ou les classes.}

\newpage 

\subsection*{\textsl{Albert Camus, de la transfiguration}, Laurent Bove}$ $\\

\subsubsection{Vérité innocente, vérité meurtrière} 
Notre ère fut celle de la violence guerrière et idéologique le siècle dernier. La violence, nous le savons, s'exerce dès que l'homme ne peut plus s'affirmer ; la voie de survie est celle de la logique et de l'efficacité. Or les héros camusiens offrent une autre fenêtre sur la façon d'être au monde. Plus que tous les autres, Meursault – l'\textsl{\'Etranger} – subit la peine indolore que la société des victimes et des juges semble reconnaitre dans son indifférence. Certes Meursault tire, aveuglé par le soleil et le reflet de la mer, « trois ou quatre coups » de pistolet sur un homme à mort. Or sa peine est aussi celle de sa froideur lors de l'enterrement de sa mère. Quelles raisons nous reste-il pour distinguer le bon du mauvais, l'espoir du désespoir ? Débarrassés des illusions matérielles et idéologiques, seulement la persévérance et la force d'exister permettent la distinction.

Il y a deux Camus : celui de \textsl{Noces} et de \textsl{l'Envers et l'Endroit}, qui s'intéresse à l'individu placé dans le monde, face à la mort des idéaux et de son corps. Mais il y a \textsl{L'Homme révolté} et \textsl{La Peste}, les \textsl{Lettres à un ami allemand} ainsi que les \textsl{Carnets}. C'est en communauté qu'évolue l'homme, et l'écrivain lutte contre le procès de l'Histoire, la domination du temps et de l'éternel présent. Comment concilié ces deux faces d'une même pensée ? Il y a le temps de la plénitude et du combat, et Laurent Bove concilie ces deux limites en une perpective spinoziste d'une révolte effective et commune. Partout, se rappeler ceci : « consentir à la vérité de la nature, à son absolue immanence, c'est en effet, dans un même geste consentir à une vérité aussi meurtrière qu'innocente. » 


\subsubsection{Plénitude de la présence}Il n'y a pas de frontière entre l'homme et le monde, de même qu'il n'y a pas de frontière entre l'enfant et la mère ; entre la mer et le ciel. Meursault reçoit la mort de sa mère, et en même temps sa propre mort, sur le mode d'une Annonciation négative. La religion de Meursault, à savoir ce qu'il reste du désir ou de l'action une fois libérées des affects passifs (\textsl{\'Ethique}, Prop. 37, Scolie 1), est celle du projet philosophique de l'oubli de soi, du dépassement des illusions et du bruit. « Vivre, c'est un peu le contraire d'exprimer. Si j'en crois les grands maîtres toscans, c'est témoigner trois fois, dans le silence, la flamme et l'immobilité. » 

Que ce qui rempli le coeur et l'esprit face à la mort et au silence ? Ce n'est ni l'adoration d'une figure christique ni l'observation des dogmes d'un quelconque contrat au monde (Goethe, Rousseau). Pourtant Camus a beaucoup de respect envers la \textit{personne} d u Christ, dont la philosophie n'a jamais le refuge dans un arrière-monde ou dans une intériorité dénaturée, mais celle qui fait de l'éternité un mensonge. En effet, Meursault poursuit dans son amour limpide de Marie, dans son observation juvénile du monde, une \textit{dépossession} et un \textit{détachement} du fantasme et de l'imagination (pervers – cf. la demande de fiancailles de Marie) — « Malheureux le pays qui a besoin d'un héros. » (\textsl{La vie de Galilée})

Meursault, une fois prisonnier, retourne à ce silence auquel on pourrait opposer le fourmillement psychologique de l'enfermement. Il partage de pathos dépouillé des figures de Piero della Francesca, la « grandeur farouche du \textsl{Christ réssussité} » qui ne demande qu'à vivre. Il sont indifférents aux yeux du spectateur et du monde. Un même mouvement accompagne le refus de la vue d'un prêtre de Meursault et de la matérialité du peintre : un déthéologisation. Ainsi les corps puissants connaissent la \textit{résurrection} :
\begin{center}
« Comme ces hommes que beaucoup de science ramène à Dieu, beaucoup d'années ont ramenés les ruines à la maison de leur mère. Aujourd'hui enfin leur passé les quitte et rien ne les distraits que cette force profonde qui les ramène au centre des choses qui tombent. » 
\end{center}
\vspace{5pt}

\subsubsection{Résistance et histoire}

Nous l'avons dit, le monde de Camus est pourtant celui du temps dirigé, de l'efficacité et de la négation nihiliste. Il entreprend avec \textsl{L'Homme révolté} la fondation d'une révolte active, dépassant la double négation hégélienne. Comment la grâce du dénuement que nous avons étudier peut éclairer cette résistante ? Notons deux perceptions du temps distinctes : selon qu'on le perçoit comme don (le temps qui passe), ou comme produit de l'activité historique (le temps qui se créer \textit{a posteriori}, \textit{sub specie aeternatis}). L'\textsl{\'Etranger} présente une domination naturelle du temps dans « cette entente amoureuse de la terre et de l'homme délivré de l'humain. » Il s'agit de concilier le don de l'actuel et l'actuel comme don.

La totalisation de l'histoire débute avec le christianisme, et se prolonge au \textsc{xx}{ème} siècle avec la matrice marxiste d'une part, sociologique (cf. P. Bourdieu) et capitaliste (cf. G. Lipovetsky) d'autre part. Le christianisme \textit{créer} un temps, un lieu, un imaginaire, duquel l'homme est séparé — à la limite pourra-t-il y pénétrer et y accrocher quelques espérances forcément déçues. Le message du Christ se voit travestit de l'idée de \textit{jugement} et de \textit{châtiment}. Voient alors le jour les philosophies du contrat (que Bove qualifie étrangement de bourgeois) et de l'histoire comme totalité : nazisme, marxisme – l'\textit{histoire meurtrière}. Pour Camus, c'est contre l'illusion de la transcendance et de la finalité que la communauté doit se dresser. Plus difficile encore, elle se doit de maintenir sa position entre révolte et révolution.
\vspace{5pt}

Nous avons vu la transfiguration est une même mesure \textit{pour} la liberté et \textit{contre} la domination du temps et de l'Histoire. La mesure est placée entre le réel et le désir, comme l'absurde est placée entre le monde et l'expérience. Pour la liberté de l'homme dans le monde d'une part, dans l'histoire d'autre part. La révolte n'est pas éloigné du prolongement politique de la transfiguration du premier Camus : « Aller jusqu'au bout, ce n'est pas seulement résister, c'est aussi se laisser aller. J'ai besoin de sentir ma personne, dans la mesure où elle est le sentiment de ce qui me dépasse. » (\textsl{Carnets I})

\newpage 

\subsection*{Science économique et économie financière}$ $\\

Sont rassemblés ci-dessous la lecture des ouvrages de Thomas Angeletti : \textsl{L'invention de l'économie française}, d'André Orléan : \textsl{L'empire de la valeur} et de Liêm Hoang-Ngoc : \textsl{Petit manuel critique des théories économiques}, accompagnés par de nombreuses pages Wikipédia\footnote{Parmi les pages les plus éclairantes se trouvent : \href{https://fr.wikipedia.org/wiki/Accords_de_Bretton_Woods}{Accords de Bretton Woods}, \href{https://fr.wikipedia.org/wiki/Histoire_de_la_dette_publique_française}{Histoire de la dette publique française}, \href{https://fr.wikipedia.org/wiki/Histoire_des_dépenses_publiques_françaises}{Histoire des dépenses publiques françaises}, \href{https://fr.wikipedia.org/wiki/Déficit_budgétaire_et_déficit_public}{Déficit budgétaire et déficit public}, \href{https://fr.wikipedia.org/wiki/Système_monétaire_international}{Système monétaire international}, \href{https://fr.wikipedia.org/wiki/Trente_Glorieuses}{Trente Glorieuses}, \href{https://fr.wikipedia.org/wiki/Monétarisme}{Monétarisme}, \href{https://fr.wikipedia.org/wiki/Marché_financier}{Marché financier}, \href{https://fr.wikipedia.org/wiki/Modèle_scandinave}{Modèle scandinave}.}. 

\subsubsection*{Perception historique de l'économie}

\paragraph{Pour comprendre la genèse, la conceptualisation de la science économique et les miracles et les crises de l'économie financière, il faut décrire en détail le processus d'autonomisation de ces deux domaines. L'économie est un champ agrégeant de multiples disciplines et son étude est un va-et-vient constant entre la politique, la sociologie, les modèles économiques et leurs mathématiques. Pour clarifier, on distinguera la \textit{science économique} de l'\textit{économie financière} (désignées respectivement \textit{economical science} et \textit{economics} en anglais).}

\paragraph{Au départ l'économie n'existe que comme processus d'échange de biens et de services. C'est le stade primitif d'une économie entre individu qui cherche à répondre à leur besoin. La \textit{monnaie} est la forme marchande de l'argent, qui permet l'existence d'un marché commun au sein d'une société d'individus coopérants selon les mêmes rapports. Orléan avance qu'au même titre que le langage, les mécanismes d'hégémonie et de concurrence font qu'une unique monnaie émerge dans une société donnée. 

À la base de ce système économique se trouve une croyance commune étonnante : la \textit{valeur} est attribuée par consensus à la monnaie et permet l'échange de bien selon des taux fixes. Nous reviendrons sur les problèmes que pose cette croyance. Implicitement au sein des sociétés primitives se trouvent aussi des moeurs communes : du don suit le contre-don, le temps dispose d'une certaine valeur, etc.}

\paragraph{Au \textsc{xviii}{ème} se développent des économies de taille nationale voir internationale. On essai de trouver ce qui fonde l'échange économique, et quelques penseurs se détachent. Chez les physiocrates, elle se situait dans la terre (conception holiste). Adam Smith le premier attribue ce qui vaut au \textit{travail}, i.e. le temps passer au champ, ou de manière équivalente, une certaine quantité de blé produite. Les facteurs limitants à la richesse sont alors le taux de production (de la technique) ; et l'étendu du marché (des champs de blé). Ricardo qui perçoit les progrès mécaniques en agronomie déplace la notion de valeur marchande dans la \textit{rareté}. À défaut de meilleur référence, il propose l'or comme moyen de fonder l'échange économique. Dans tous les cas, il s'agit de théories substantielles de la valeur. William Marx remarque le premier la primauté du rôle des échanges sociaux dans l'économie marchande et \textit{sépare} la marchandise de la production de cette marchandise.}

\paragraph{Un siècle plus tard l'économie de marché est assise dans la plupart des pays occidentaux. Les \textit{actions} naissent à Toulouse au \textsc{xiii}{ème} siècle.}












\end{document}